\documentclass[11pt]{article}
\usepackage[a4paper,margin=24mm]{geometry}
\usepackage[T1]{fontenc}
\usepackage{lmodern,amsmath,amssymb,amsthm,mathtools,booktabs,microtype}
\usepackage[colorlinks=true,linkcolor=blue,urlcolor=blue,citecolor=blue]{hyperref}
\newtheorem{proposition}{Proposition}
\newcommand{\eps}{\epsilon}
\newcommand{\ket}[1]{|#1\rangle}
\newcommand{\B}{\widehat B}
\newcommand{\PP}{\widehat P}
\newcommand{\Bt}{B_{\rm tr}}
\newcommand{\Pois}{\operatorname{Pois}}
\title{A controlled thermal window for chord generators\\
\large Endpoint preparation and the Hamiltonian boost}
\author{Alan Day \and Chengheng\thanks{AI research assistant (OpenAI).}}
\date{8 November 2026\quad Local research draft}
\begin{document}\maketitle
\begin{abstract}
The zero-chord matter seed fails a direct identification of the
Lin--Stanford length-preserving generators with the microscopic boost.
We study their exact action on an endpoint Euclidean preparation.
A positive-path comparison with the ordinary oscillator gives a
sufficient joint window in which the inverse-temperature parameter
vanishes, the chord number diverges, and the normalized heat vector
approaches a moving coherent vector. In this window the actions of
the rescaled chord boost and the negative Hamiltonian boost agree
in norm on the prepared vector. Lowest-weight and second-moment
limits follow. The preparation becomes orthogonal to its original
seed, so the result supplies a controlled moving family, rather
than an identification of the original selected state.
\end{abstract}

\section{Conventions and the question being tested}
First take the established fixed-$\lambda$ double-scaled,
half-filled limit of the project's complex Gaussian CAR model.
The subsequent limit in this note is
\[
 q=e^{-\eps},\qquad \eps=4\lambda,\qquad
 K=\frac{H}{2J_s}=\frac{h}{\sqrt\eps},\qquad \hbar=1.
\]
In Lin--Stanford (LS) units their deformation parameter is our
$\eps$, and their Hamiltonian is our $K$ \cite[(18)]{LS}.
We work in the tracial one-insertion chord module of fixed weight
$d>0$, with Hamiltonian--matter crossing $r=q^d$.
For the established actual charged curvature channel, $d=2$;
we use that module, without adding an independent random probe.
The full selected density is not inserted in this note.

Let $\ket{l,k}=\ket{d;l,k}$, and set $[n]=(1-q^n)/(1-q)$.
LS creation and annihilation operators give
\begin{align}
 a_L\ket{l,k}&=[l]\ket{l-1,k}+rq^l[k]\ket{l,k-1},\\
 a_R\ket{l,k}&=[k]\ket{l,k-1}+rq^k[l]\ket{l-1,k}.
\end{align}
Creation adds a chord on the indicated side. The total level is
$\bar n=l+k+d$ and different total levels are orthogonal.
We compare $\Bt=(K_L-K_R)/2$ with the LS combinations of
$J_{ij}=a_i^\dagger a_j-[\bar n]$.
The new estimates below use this existing representation and do
not constitute a new chord algebra or an independent bulk dual.

\newpage
\section{A sufficient joint window from positive paths}
The endpoint choice of the LS preparation \cite[(141)]{LS} is
\begin{equation}
 \Psi=Z_q^{-1/2}e^{-\beta K_L/2}\ket{0,0},\qquad
 t=\frac{\beta}{2\sqrt\eps},\quad m=t^2,\quad
 Z_q=\|e^{-tT_q}e_0\|^2 .\label{heat}
\end{equation}
Here $e_n=\ket{n,0}/\sqrt{[n]!}$ is orthonormal and
$T_q=\sqrt\eps K_L$ has off-diagonal entries $\sqrt{[n]}$.
This endpoint subspace is invariant under $K_L$, but not under
the other generators. After a parity transformation the heat
vector has positive coefficients $x=e^{tT_q}e_0$.
The comparison vector $y=e^{tT_1}e_0$ satisfies
\[
 0\le x_n\le y_n,\qquad \|y\|^2=e^{2m},\qquad
 \chi_n=(-1)^n e^{-m/2}\frac{m^{n/2}}{\sqrt{n!}}.
\]
Thus $|\chi_n|^2$ is a Poisson law of mean $m$.

\begin{proposition}\label{comparison}
For $\delta=2\eps m^2<1$,
\begin{align}
 1-\delta&\le Z_q/e^{2m}\le1,\label{partition}\\
 \|\Psi-\chi\|^2&\le2(1-\sqrt{1-\delta})\le2\delta,\label{distance}\\
 |\Psi_n|^2&\le\frac{e^{-m}m^n}{(1-\delta)n!},\qquad
 |\Psi_0|\le\frac{e^{-m/2}}{\sqrt{1-\delta}}.\label{domination}
\end{align}
The identification in \eqref{distance} uses the orthonormal
endpoint basis for each $q$, not the raw chord basis.
\end{proposition}
\begin{proof}
The $2j$-th vacuum moment of $T_q$ sums $q^{\mathrm{cr}}$
over pairings. Each has at most $j(j-1)/2$ crossings, so its
deficit from the $q=1$ weight is at most
$\eps j(j-1)/2$. In the normalized series for
$\langle e_0,e^{2tT_1}e_0\rangle$, $j$ has law $\Pois(2m)$.
Its factorial moment gives total deficit at most $2\eps m^2$.
Positivity gives $\langle x,y\rangle\ge\|x\|^2$, proving
\eqref{distance}. The componentwise bound and \eqref{partition}
give \eqref{domination}; for the last bound use $y_0=e^{m/2}$.
\end{proof}

The sufficient, nonempty window is
\begin{equation}
 m\longrightarrow\infty,\quad \eps m^2\longrightarrow0,
 \qquad\text{equivalently}\qquad
 \sqrt\eps\ll\beta\ll\eps^{1/4}.\label{window}
\end{equation}
In particular,
$\mathbb E_\Psi(N/m-1)^2\le[(1-\delta)m]^{-1}$.
Thus $N/m\to1$, $\eps N\to0$, and the seed overlap tends to
zero. This is a sufficient window, not an optimal one. The
fixed-temperature LS saddle motivates it but does not by itself
prove this joint limiting statement.

\newpage
\section{Exact generator action at the endpoint}
In the sector $l+k=n$, put $c_n=q^{(n+d)/2}$ and use
the exact LS definitions \cite[(48), (147)]{LS}:
\begin{align}
 E&=-\frac{J_{LL}+J_{RR}}{2c_n},&
 B&=\frac{J_{LL}-J_{RR}}{2c_n},&
 P&=\frac{i}{2}(J_{LR}-J_{RL}),\\
 \B&=\frac{B}{\sqrt{1-c_n^2}},&
 \PP&=\frac{P}{\sqrt{1-c_n^2}}.&&
\end{align}
The rescaling is sectorwise. All these operators preserve $n$.
For $n\ge1$, define
\[
 u_n=\frac{\ket{n-1,1}}{\sqrt{[n]!}},\qquad
 \langle e_n,u_n\rangle=r,\qquad
 \|u_n\|^2=r^2+\frac{1-r^2}{[n]},\qquad
 f_n=\frac{u_n-re_n}{\sqrt{(1-r^2)/[n]}}.
\]
The Gram identities follow by moving a creation operator to
the bra and applying its annihilation adjoint. They make $e_n,f_n$
orthonormal, avoiding subtraction of nearly parallel raw vectors
in the numerical implementation. Write $g_n=\sqrt{[n](1-r^2)}$.
Direct application of the annihilators gives
\begin{align}
 \B e_n&=\frac{[n](1-r^2)e_n-rg_n f_n}
                 {2c_n\sqrt{1-c_n^2}},\label{Bexact}\\
 \PP e_n&=-\frac{ig_n f_n}{2\sqrt{1-c_n^2}},\label{Pexact}\\
 E e_n&=\frac{[d]-(1-r)^2[n]/2}{c_n}e_n
                 -\frac{rg_n}{2c_n}f_n.\label{Eexact}
\end{align}
At $n=0$ the $f_n$ terms are absent, $B=P=0$ and
$E=[d]/q^{d/2}$, recovering the previously recorded seed obstruction.

For an arbitrary endpoint vector $\psi=\sum_n\psi_ne_n$,
the level-$n$ components of the physical boost are instead
\begin{align}
 (\Bt\psi)_{e,n}
 &=\frac{1-r}{2\sqrt\eps}
       \left(\sqrt{[n]}\,\psi_{n-1}
                    +\sqrt{[n+1]}\,\psi_{n+1}\right),\label{physicale}\\
 (\Bt\psi)_{f,n}
 &=-\frac{\sqrt{1-r^2}}{2\sqrt\eps}\,\psi_{n-1}.\label{physicalf}
\end{align}
Missing coefficients are zero. Equations \eqref{physicale}--\eqref{physicalf}
retain both lowering and raising, including the first output
level beyond a numerical cutoff. The comparison is between
operators in the same chord Hilbert space, rather than between
two independently fitted correlation functions.

\newpage
\section{Matching on the moving preparation}
\begin{proposition}\label{matching}
For each fixed $d>0$, in \eqref{window},
\begin{align}
 \|(\B+\Bt)\Psi\|&\longrightarrow0,&
 \|(E-d)\Psi\|&\longrightarrow0,&
 \|(\B+i\PP)\Psi\|&\longrightarrow0,\label{strong}\\
 \langle\B\rangle_\Psi,\ \langle\Bt\rangle_\Psi&\longrightarrow0,&
 \langle\B^2\rangle_\Psi,\ \langle\PP^2\rangle_\Psi,
 \langle\Bt^2\rangle_\Psi&\longrightarrow d/2.\label{moments}
\end{align}
Also $\langle[\B,\PP]\rangle_\Psi\to id$. These are
one-vector and weak-matrix-element limits; no convergence of
commutators in norm or of the full operators is asserted.
\end{proposition}
\begin{proof}
First work on the comparison vector $\chi$.
On levels with $n\to\infty$ and $\eps n\to0$,
\[
 [n]/n\to1,\quad(1-r^2)/\eps\to2d,\quad
 (1-c_n^2)/(\eps(n+d))\to1.
\]
Equations \eqref{Bexact}--\eqref{Eexact} therefore give,
uniformly in the Poisson bulk as $n/m\to1$,
\[
 \B e_n=-\sqrt{d/2}\,f_n+o(1),\qquad
 \PP e_n=-i\sqrt{d/2}\,f_n+o(1),\qquad E e_n=de_n+o(1).
\]
Here vector remainders are in the level's orthonormal coordinates.
The coherent coefficients satisfy exactly
$\chi_{n-1}=-\sqrt{n/m}\,\chi_n$ and
$\chi_{n+1}=-\sqrt{m/(n+1)}\,\chi_n$.
Consequently \eqref{physicalf} approaches
$+\sqrt{d/2}\,\chi_n f_n$, while \eqref{physicale} vanishes:
its bulk size is $O_d(\sqrt{\eps m})$.
This proves the plus sign in \eqref{strong}.

To justify the bulk restriction and transfer from $\chi$ to
$\Psi$, one needs more than norm closeness alone. The squared
endpoint coefficients of $\B,\PP,E$, and the two shifted
coefficients of $\Bt$, obey a common envelope
\begin{equation}
 C_d(1+\eps(n+1))^2 e^{2\eps(n+d+1)},\qquad0<\eps\le1.\label{envelope}
\end{equation}
For example, use $[n]\le n$, $1-r^2\le2d\eps$ and
$1-e^{-z}\ge ze^{-z}$; the diagonal part of $\B$ squared
is at most $d^2\eps n e^{2\eps(n+d)}$, and its $f_n$
part squared at most $(d/2)e^{2\eps(n+d)}$.
The other coefficients follow from the same inequalities.
On $m/2\le n\le2m$ the envelope is bounded uniformly.
On its complement, the Poisson domination \eqref{domination}
and exponential tilting of $\Pois(m)$ make the weighted tail
tend to zero. Indeed, tilting by $e^{2\eps n}$ gives mean
$me^{2\eps}$ and prefactor $e^{m(e^{2\eps}-1)}\to1$;
the fixed polynomial weights do not overcome the exponential
concentration. Shifts by one level have the same property.
Equation \eqref{distance} now transfers each graph norm.
The moments follow by inner products; specifically
$\langle[\B,\PP]\rangle=2i\operatorname{Im}
\langle\B\Psi,\PP\Psi\rangle\to id$.
\end{proof}

\newpage
\section{Finite checks, truncations, and numerical corrections}
The exact rational checks use $q=1/4,4/9$, $d=1,2,3$ and
levels $0\le n\le7$. They test the full chord Gram symmetry
and creation/annihilation adjoint relation, endpoint Gram and
generator formulas, and both boost degree shifts. Independent
pairing enumeration through order twelve checks the positive
vacuum-moment input. There are 4,152 exact equalities in total.

For the thermal vectors we use $m=8,32,128,512$,
$\eps=m^{-3}$, $\beta=2/m$, so $\delta=2/m$.
The coarse Jacobi cutoffs are $122,226,530,1522$; each is
increased by 64. The scaled positive Taylor vector is
$e^{-m}\sum_{j=0}^{L}(tT_{q,M})^je_0/j!$.
The actual Taylor orders are $274,570,1451,3628$.

Here are bounds that avoid numerical underflow. Let
$\pi_\mu(k)=\Pr\{\Pois(\mu)\ge k\}$.
Positive paths and their maximum level give, in exact arithmetic,
\begin{align}
 \|\Psi-\Psi_M\|^2&\le
           \frac{2\pi_{2m}(M+1)}{1-\delta},\label{cutoff}\\
 \left\|e^{-m}e^{tT_{q,M}}e_0-
 e^{-m}\sum_{j=0}^{L}\frac{(tT_{q,M})^j}{j!}e_0\right\|^2
 &\le\pi_{2m}\bigl(\lceil(L+1)/2\rceil\bigr).\label{Taylor}
\end{align}
For \eqref{cutoff}, a vacuum path visiting level $M+1$
has at least $M+1$ up steps. For \eqref{Taylor}, positivity
bounds the squared remainder by the difference of squared
norms. Missing pairs of Taylor powers have even total degree
$2j\ge L+1$. In both cases compare to the $q=1$ moments.
We evaluate the tails using
$\log\pi_\mu(k)\le-\mu+k(1+\log(\mu/k))$, $k>\mu$.
The largest vector bounds are $1.68\times10^{-23}$ for the
Jacobi cutoff and $8.41\times10^{-39}$ for the scaled Taylor tail.

\begin{center}
\begin{tabular}{rrrrr}
\toprule
$m$ & $\langle N\rangle$ & $\|(\B+\Bt)\Psi\|$
 & $\|(E-2)\Psi\|$ & $\langle\B^2\rangle$\\
\midrule
8   & 7.897924 & 0.171789 & 0.125648 & 0.837580\\
32  & 31.973991 & 0.087559 & 0.031260 & 0.943929\\
128 & 127.993490 & 0.044095 & 0.007813 & 0.984796\\
512 & 511.998372 & 0.022085 & 0.001953 & 0.996120\\
\bottomrule
\end{tabular}
\end{center}
These entries use the actual channel $d=2$; the same driver
also checks $d=1,3$. At $m=512$, $\langle\Bt^2\rangle=1.000015$,
whereas $\|(\B-\Bt)\Psi\|=1.997945$, confirming the sign.
The cancellation ratio
$(\|\B\Psi\|+\|\Bt\Psi\|)/\|(\B+\Bt)\Psi\|$ is at most
$90.48$ for this channel. It is recorded rather than suppressed.

\newpage
\section{What the comparison establishes and what comes next}
The single-step matching and the limiting lowest-weight
relations provide a concrete route beyond the old zero-chord
test. They do not leave the preparation fixed. In fact,
\eqref{domination} proves $\langle e_0,\Psi\rangle\to0$;
at $m=512$ the overlap is about $6.62\times10^{-112}$.
The same bulk/tail argument, applied to $T_q/\sqrt m$, gives
\[
 \|(T_q/\sqrt m+2)\Psi\|\longrightarrow0.
\]
Thus the thermal vector moves far in rms Hamiltonian units even
though $\beta\to0$. A claim that this regulator simply restores
the original trace preparation would be incorrect.

The original full selected preparation has
\[
 \rho=1-\sqrt q(1+\lambda)+\sqrt q(h-\sqrt\lambda U)^2.
\]
Its $U$ dependence is not represented by the endpoint heat
vector above. A useful next step is an explicit moving-space
identification, followed by mixed generator matrix elements
with this density retained. A displacement on the endpoint
oscillator alone has not yet been extended to the entire matter
module. Products of generators, commutator norms on descendants,
and the selected mixed Ward identities remain to be established.
Neither $d=2$ nor $E\Psi\simeq2\Psi$ fixes a Poincare mass;
positive null translations and their multiparticle composition
still require a separate construction. The selected four-bilocal
and finite-$G$ dynamics are also outside this result.

\paragraph{Numerical record.}
Long calls to \texttt{expm\_multiply} failed the initial
precision check; splitting and normalizing reduced, but did not
eliminate, the discrepancy. The largest split-method vector
difference is $3.45\times10^{-9}$, and this method is retained
only as a diagnostic. The positive Taylor calculation was
checked at 50 decimal digits: the largest relative discrepancy
among the eight recorded $d=2$ observables is
$3.16\times10^{-11}$, in the small coherent-distance quantity.
The initial \texttt{longdouble} comparison did not increase
precision on this Windows host; the supplemental decimal check
explicitly corrects that description. An underflowed zero in
the original Taylor-bound field is likewise superseded by
\eqref{Taylor} and the separate logarithmic tail audit.
The bounds are for exact arithmetic, not interval certificates
for all floating-point operations. Failed diagnostics and both
corrections are preserved alongside the sealed scientific outputs.

\paragraph{Predecessors and attribution.}
The q-oscillator, one-matter representation, Euclidean
preparation and rescaled generators are LS inputs. Central-spectrum
coherent behavior and thermal chord wavefunctions are already
studied by Xu \cite[Secs.~2.3, 5]{Xu}; q-coherent preparations
are also discussed in \cite[Sec.~3.2.1]{FJKP}. Our ordinary
exponential heat vector is not defined to be a q-annihilation
eigenstate. This note supplies a sufficient norm-controlled
joint window and an endpoint application of the existing
generators. We do not claim priority for coherent-state methods
or for emergent $\mathfrak{sl}_2$ symmetry.

\begin{thebibliography}{9}
\bibitem{LS} H.~W. Lin and D. Stanford,
\emph{A symmetry algebra in double-scaled SYK},
\href{https://arxiv.org/abs/2307.15725}{arXiv:2307.15725};
formula numbering from the v1 HTML.
\bibitem{Xu} J. Xu, \emph{On Chord Dynamics and Complexity Growth
in Double-Scaled SYK},
\href{https://arxiv.org/abs/2411.04251v3}{arXiv:2411.04251v3}.
\bibitem{FJKP} Y. Fu, H.-S. Jeong, K.-Y. Kim and J.~F. Pedraza,
\emph{Toward Krylov-based holography in double-scaled SYK},
\href{https://arxiv.org/abs/2510.22658v1}{arXiv:2510.22658v1}.
\end{thebibliography}
\end{document}
